\section{Witnesses concerning the first angelic sin}

\begin{longtable}{>{\raggedright\arraybackslash}p{.21\linewidth}>{\raggedright\arraybackslash}p{.43\linewidth}>{\raggedright\arraybackslash}p{.26\linewidth}}
\textbf{Witness and locus} & \textbf{Determination} & \textbf{Question left distinct}\\\hline\endhead
Augustine, \emph{City of God} XI.13--15 & The devil's nature was God's good work; his abandonment of truth belongs to his will. Pride explains the refusal of rightful subjection. & The account of the angels' original condition and the timing of their assurance of perseverance.\\
John of Damascus, \emph{Exposition} II.4 & A good angel assigned to the earthly realm voluntarily rebelled; other angels followed him. & This assignment does not establish Thomas's preferred account of the rebel's original rank.\\
Bonaventure, \emph{Breviloquium} II.7 & Lucifer presumptuously rested in his private excellence and sought superiority over others; hatred and envy toward humanity follow the fall. & Created excellence itself remains a good, distinct from its proud appropriation.\\
Thomas, \ST{I, q.~63, aa.~2--3} & Pride is the first sin; envy follows. The rebel desires likeness to God through an inordinate mode of attaining beatitude. & Desire for divine likeness does not mean believing that a creature could become God by nature.\\
Thomas, \ST{I, q.~63, a.~7} & The highest angel's fall is judged more probable; the Damascene account is not rejected as contrary to faith. & The chief rebel's rank is an opinion distinct from the doctrine of the voluntary fall.\\
\end{longtable}

\noindent\textit{Sources:} Augustine, \emph{City of God}, trans.
Marcus Dods, book XI, chapters 13--15,
\url{https://www.newadvent.org/fathers/120111.htm}; the Damascene,
Bonaventure, and Thomas witnesses are identified in the preceding chapters.
The will's rejection of God and its ensuing envy stand in a logical
sequence: envy proceeds from the proud attachment to a private good.
